\begin{textsample}{POS Dim 1 – se – Score 43.00 – se/se\_inf\_e\_ef\_2.txt}  \label{ex:f1_pos_112}
1.
TEXTO INTRODUTÓRIO : TRAVESSIAS NO ENSINAR E APRENDER \textbf{Um} projeto pedagógico escolar deve ser \textbf{um} convite à construção da cidadania, \textbf{um} convite à promoção e à afirmação de valores \textbf{que} pressuponha a \textbf{dignidade} humana.
E esse projeto deve materializar -se num currículo capaz de colocar estudante e professor em diálogo fecundo, com a experiência histórica da comunidade na qual estão inseridos e, a partir dela, com os mais diversos tempos e espaços.
O ensinar e o aprender em todas as escolas, em cada sala e noutros espaços \textbf{que} o currículo venha a articular, devem cumprir os propósitos éticos, estéticos e políticos dispostos na legislação educacional brasileira, não esquecendo àqueles fixados na Carta Constitucional.
A ciência, seus saberes \textbf{sistematizados}, as linguagens e o \textbf{método} filosófico devem ser as referências fundamentais desse projeto curricular, sem \textbf{que} com isso se negue todo e qualquer conhecimento \textbf{que} constitua a experiência humana, pelo contrário, é com todos eles \textbf{que} os saberes de referência devem dialogar, possibilitando, num exercício crítico, a \textbf{problematização} das realidades.
A aprendizagem \textbf{que} emergir desse fazer educacional deverá ter como foco o pensar e a (re)significação dos saberes e fazeres, dos modos de intervenção e de produção das existências individuais, compreendendo \textbf{que} são elas também existências coletivas. \textbf{Um} pensar \textbf{que} se confunda com o experimentar, \textbf{que} se faça sinônimo do sentir.
Assim sendo, serão as linguagens e as artes, simultaneamente, saberes a serem desenvolvidos e suportes para \textbf{que} a \textbf{tarefa} de aprender e ensinar as ciências desenvolva -se.
O \textbf{sujeito} \textbf{que} aprende e ensina é o \textbf{sujeito} \textbf{que} se afeta pelas coisas do mundo, pelo outro e, do mesmo modo, afeta ao outro.
Cada escola básica está inserida numa comunidade e, por integrar \textbf{uma} rede, ela deve falar com e para a sua cidade.
Como a cidade, na arquitetura política da República brasileira, está contida numa unidade federativa e, a partir dela, num Estado-nação, o currículo escolar não deve negligenciar a diversidade sociocultural, étnica e geo-histórica dessas outras \textbf{territorialidades}.
Por isso, as diretrizes curriculares nacionais devem ser a base legal e \textbf{conceitual} desse projeto.
É fundamental \textbf{que} o coletivo escolar construa e compartilhe de \textbf{um} entendimento acerca do currículo.
Pode ter ele \textbf{uma} pluralidade imensa de significados.
Perpassa diferentes espaços e \textbf{instâncias}.
É elemento propulsor e regulador.
Carrega intencionalidades \textbf{que} variam, de acordo, com as perspectivas adotadas.
Para a Educação Básica, no estado de Sergipe, propõe -se \textbf{que} a escola pense o currículo como \textbf{uma} rede de saberes, caminhos e práticas, \textbf{que} articulam todas as dimensões do projeto escolar ao propósito de fazer de cada criança e adolescente \textbf{um} cidadão comprometido eticamente com o cuidado de si e do outro, com o exercício crítico e responsável da cidadania, com sua cidade, com a defesa da diversidade étnica e cultural brasileiras, com o patrimônio \textbf{sócio-histórico} e ambiental do estado e do país, com a democracia, com \textbf{um} desenvolvimento social, inclusivo e sustentável.
Para a escola \textbf{que} se pretende \textbf{contemporânea} e compromissada com esse conjunto de coisas, não há outra opção senão pensar o \textbf{sujeito} \textbf{que} aprende em sua \textbf{integralidade}.
O aprender é \textbf{um} exercício \textbf{que} se \textbf{opera} num corpo.
Num corpo moldado pelas relações com o mundo. \textbf{Um} corpo \textbf{que} é afeto de \textbf{cognição} e \textbf{que}, em razão do tempo e do espaço, dos valores e costumes erguidos sob diversas cronologias e \textbf{territorialidades}, acaba \textbf{que} sentindo/pensando e \textbf{operando} sobre o mundo de modos diferentes.
O projeto escolar precisa ser integral, antes na \textbf{abordagem} \textbf{que} no tempo, porque os saberes socialmente produzidos e \textbf{historicamente} especializados \textbf{dizem}, em última \textbf{instância}, do “ \textbf{homem} ” e das coisas do mundo \textbf{que} ele inventou para si.
São, portanto, modos distintos de falar sobre as mesmas coisas.
Perspectivas variadas \textbf{que} possibilitem a compreensão da \textbf{integralidade}, seja do \textbf{sujeito} \textbf{que} aprende, seja do objeto \textbf{que} se quer conhecer.
Não se trata de ensinar e aprender tudo.
É \textbf{tarefa} impossível.
O desafio é menos de quantidade e mais do modo como se concebe o \textbf{sujeito} \textbf{que} aprende e o objeto da aprendizagem.
Cada criança e adolescente é \textbf{um} ser único, mas é também \textbf{um} ser múltiplo, porque é sócioemocional, cognoscente, cultural e corpóreo.
Não se abandona qualquer dessas perspectivas quando se aprende.
A \textbf{abordagem} acerca do objeto do conhecimento e o tempo/espaço em \textbf{que} ela se dá deve, sempre, considerar todos esses aspectos, porque de todos eles \textbf{depende}, simbioticamente, a aprendizagem.
A escola \textbf{que} assumir o \textbf{compromisso} com essa perspectiva de ensino integral será, \textbf{certamente}, \textbf{uma} escola inclusiva.
Porque reconhecerá e promoverá a alteridade ; porque terá como foco a \textbf{integralidade} e, assim sendo, compreenderá os tempos e limites singulares da \textbf{tarefa} de aprender.
E será inclusiva, sobretudo, porque possibilitará a cada estudante o convívio com o diverso, com a multiplicidade abrigada em si.
Com o outro \textbf{que}, em razão de certas diferenças, se não estará na escola para aprender sobre as “ ciências ”, estará, sobretudo, para ensinar e aprender a exercer sua \textbf{humanidade}.
Do mesmo modo, é \textbf{imprescindível} \textbf{que}, na escola, todos aqueles \textbf{que} a fazem compartilhem algum conceito de ensinar e aprender ; \textbf{que} reconheçam e afirmem \textbf{que} a escola possui \textbf{uma} função político-social.
Não compete a ela, solitariamente, a construção da cidade, portanto, da cidadania.
Mas tem ela, nessa complexa \textbf{tarefa}, \textbf{um} papel fundamental.
Ensinar e aprender são \textbf{faces} de \textbf{uma} mesma relação com o mundo, com o outro, com os saberes \textbf{que} são produzidos nas interações sociais.
São facetas dos processos de apropriação dos conhecimentos \textbf{que} se dão nessas mesmas interações.
Aprende -se e ensina -se no decorrer da vida.
Nesse processo, torna -se humano.
Ensinar e aprender são, assim, \textbf{fundamento} da condição humana.
Aprender é experimentar, é se apropriar, é transformar, é se deslocar e (re)significar o mundo.
Nessa \textbf{tarefa}, cada \textbf{um} faz a si mesmo na medida em \textbf{que} (re)faz as coisas do mundo \textbf{que} lhe circunda.
Sendo a aprendizagem \textbf{um} fenômeno eminentemente social, implícito ao universo das linguagens, o aprender traz a dimensão do ensinar. \textbf{Um} território adequado ao desenvolvimento da aprendizagem, especialmente para crianças e adolescentes, deve ser \textbf{atravessado} de afeto, de cuidado, de \textbf{estímulo} sensório-motor, portanto, de \textbf{estímulo} à imaginação.
A afecção cumpre papel fundamental na produção da memória, sendo essa, por sua vez, indispensável ao desenvolvimento da aprendizagem.
A plasticidade desses processos, na escola, não deve dispensar o planejamento e deverá sempre considerar \textbf{que} o percurso previsto poderá ser refeito, inclusive, ao tempo em \textbf{que} se caminha.
E na caminhada, a memória, a experiência acumulada devem ser lugares para visitas, nunca moradas.
O ensinar e o aprender \textbf{aqui} propostos não devem \textbf{cristalizar} os modos de fazer, de estar, de pensar.
A aprendizagem é \textbf{um} processo profundamente marcado por sua superfície sociocultural, mas não se deve esquecer \textbf{que} ela é, também, \textbf{uma} operação bioquímica, \textbf{dialeticamente} articulada às socializações.
Faz -se efetivamente no corpo.
Num corpo \textbf{que} sente o mundo e se afeta por ele, produzindo como resposta, permanentemente, \textbf{uma} memória.
O corpo, \textbf{que} foi \textbf{historicamente} humanizado, é produto de \textbf{um} tempo evolutivo de longuíssima duração.
Traz \textbf{um} sistema nervoso comandado por \textbf{um} cérebro, \textbf{que} é o mais complexo e sofisticado dentre aqueles observados nas formas de vida \textbf{conhecidas}.
Compreender os \textbf{fundamentos} neuropsíquicos da aprendizagem é \textbf{tarefa} incontornável para aqueles \textbf{que} fazem a educação escolar e pretendem \textbf{que} ela seja integral e inclusiva.
Na escola, o ensinar e o aprender devem adquirir profunda intencionalidade.
Estarão \textbf{demarcados} por tempos, espaços e propósitos específicos.
Sem a ambição de dar conta de todas as aprendizagens e de todos os ensinamentos, caberá ao coletivo escolar fazer escolhas, sempre referenciado pelas diretrizes educacionais brasileiras e pelas diretrizes da Rede à qual a unidade pertença e sempre considerando suas específicas necessidades.
Escolher caminhos, definir programas, estabelecer objetivos e metas, construir rotinas \textbf{que} façam de cada sala de aula, da escola como \textbf{um} todo, \textbf{um} fecundo ambiente de aprendizagens.
Nessa segunda década do \textbf{século} 21, os sergipanos se encontram diante de desafios \textbf{que} têm a dimensão do seu território e outros tantos \textbf{que} compartilham com os demais brasileiros.
E, nesses dias em \textbf{que} todo o planeta se \textbf{configura} como \textbf{uma} grande aldeia, não se deve furtar -se de pensar as incitações da Odisseia humana.
A escola deve estar compromissada com todas essas provocações.
São muitas as territorializações \textbf{que} lhe perpassam.
O \textbf{sujeito}, o universo virtual, o bairro, a cidade, o estado, o país, a mundialização.
O Brasil é \textbf{um} país de dimensões continentais.
De longa data, reconhecido pela exuberância e diversidade de sua fauna, flora e riqueza hídrica.
É \textbf{uma} nação etnicamente mestiça, mas, ainda assim, marcada por processos históricos de \textbf{exclusão} nos quais sobressaem questões étnicas.
Os mesmos projetos de produção social da riqueza \textbf{que} se sucederam, desde a ocupação colonial, engendraram e \textbf{cristalizaram} esquemas socioculturais de \textbf{exclusão} e, ao mesmo tempo, empreenderam, quase sempre, formas insustentáveis de exploração dos seus recursos naturais.
No \textbf{século} XX, na esteira desses movimentos, o país \textbf{viveu} \textbf{uma} das maiores marchas migratórias da história, quando no curso de quatro décadas ele deixou de ser predominantemente campesino para tornar -se hegemonicamente urbano.
Desse feito, resultaram muitos dos desafios \textbf{que} as cidades e o “ campo ” brasileiros enfrentam em pleno o \textbf{século} XXI.
O planejamento urbano, o saneamento básico, a universalização da moradia digna, a mobilidade, a produção e acondicionamento dos resíduos e as violências são apenas alguns deles.
Esse quadro impõe ao projeto escolar brasileiro, como a poucos no mundo, \textbf{um} \textbf{compromisso} inadiável com o princípio da sustentabilidade, com a democracia e com a \textbf{equidade} social.
Não se trata de apregoar \textbf{um} “ ambientalismo ” ingênuo, pois nossa presença no mundo sempre resultará em algum impacto.
Não se trata também de imputar a \textbf{um} único indivíduo as transformações pelas quais \textbf{necessitam} passar o estado e sociedade brasileiros.
Cabe -nos construir \textbf{uma} cultura social \textbf{que} problematize o padrão de produção e de consumo vigentes, e \textbf{que} proponha alternativas para a superação de muitos dos desafios \textbf{que} estão postos. \textbf{Operar} tais mudanças é \textbf{tarefa} de \textbf{uma} sociedade inteira, mas debater criticamente e comprometer os \textbf{sujeitos} com essas realidades são funções precípua da escola básica.
Não é coincidência \textbf{que} as sociedades \textbf{que} mais avançaram no cuidado com seu patrimônio ambiental são também as \textbf{que} mais radicalizaram, na \textbf{direção} da \textbf{equidade} social, seus regimes democráticos.
É fundamental \textbf{que} o aprendizado dos saberes das ciências, mobilize a criatividade das nossas crianças e jovens e credencie -os para, sob o crivo de \textbf{uma} ética, aplicar esses conhecimentos, construindo respostas para os desafios da vida, da cidade e do país, cujo objetivo maior é assegurar a \textbf{dignidade} de cada pessoa para garantir \textbf{uma} relação de equilíbrio com o mundo.
Como se materializa \textbf{um} projeto de comunidade sustentável, humana e criativa?
Como fazer da urbe \textbf{que} lhe circunda e das suas vivências \textbf{um} território inclusivo? \textbf{Que} contribuições, nesse sentido, pode cada município oferecer aos sergipanos e aos demais brasileiros?
Quais recados cada cidade poderá enviar ao mundo?
Quanto dessa \textbf{tarefa} compete ao projeto escolar?
Os caminhos \textbf{que} cada sala de aula, \textbf{que} todas as escolas traçarem nessas \textbf{direções} denotarão o \textbf{compromisso} político-social dessa instituição, sendo esse, portanto, o \textbf{compromisso} dos seus coletivos.
Não há conhecimento \textbf{neutro}.
Nunca existiu.
Toda ação educativa pressupõe \textbf{uma} posição ético-política no mundo.
Cada conteúdo selecionado, assim como cada itinerário formativo proposto, estará \textbf{atravessado} de intenções.
O \textbf{que} \textbf{exige} \textbf{que} diante de qualquer deles se deva sempre ter \textbf{uma} posição crítica.
É falsa a ideia de \textbf{que} existe \textbf{uma} educação de virtudes, supostamente de obrigação da família e \textbf{uma} educação de conteúdos escolares.

% matched lemmas: abordagem, aqui, atravessar, certamente, cognição, compromisso, conceitual, configurar, conhecido, contemporâneo, cristalizar, demarcar, depender, dialeticamente, dignidade, direção, dizer, equidade, estímulo, exclusão, exigir, face, fundamento, historicamente, homem, humanidade, imprescindível, instância, integralidade, método, necessitar, neutro, operar, problematização, que, sistematizar, sujeito, século, sócio-histórico, tarefa, territorialidade, um, viver
\end{textsample}
